% R008 — Guide to Cultivating Complex Analysis (Bahasa Indonesia)
% Reader-facing translation unit: source/ca-v1.9/ca.tex lines 734–759.
% Source release: v1.9 (commit a4d25d9ba37a486a5a020219eb8bd4e6602fd14c).
% Translation provenance: OpenAI Codex gpt-5.6-sol, Ultra, acting on the user's request.
% Preserve source labels, glossary hooks, mathematical notation, and structure.

\subsection{Bilangan kompleks sebagai bidang} 

Anda tentu pernah melihat medan bilangan kompleks,
tetapi mari kita tinjau kembali definisinya.
Sebagai suatu himpunan, misalkan \emph{\myindex{medan bilangan kompleks}}
\glsadd{not:C}$\C$ adalah himpunan $\R^2 = \R \times \R$.
Himpunan ini merupakan sebuah bidang, sehingga kita menyebutnya
\emph{\myindex{bidang kompleks}}%
\footnote{Meskipun ada seorang matematikawan aneh di luar sana yang
menganggap bahwa \emph{bidang kompleks} adalah $\C^2 = \C \times \C$.
Jika Anda mendengar seseorang mengatakan demikian, pukul kepalanya dengan
sopan untuk saya.}.
Agar menjadi medan, kita mendefinisikan penjumlahan dan perkalian:
\begin{align*}
(a,b) + (c,d) & \overset{\text{def}}{=} (a+c,b+d) , \\
(a,b) (c,d) & \overset{\text{def}}{=} (ac-bd,bc+da) .
\end{align*}

\begin{exbox}
\begin{exercise}%[\neededexmark]
Periksa bahwa $\C$ adalah medan, dengan identitas aditif
$0=(0,0)$ dan identitas multiplikatif $1=(1,0)$.
Artinya, $\C$ merupakan grup Abelian terhadap penjumlahan,
bilangan kompleks tak nol merupakan grup Abelian terhadap perkalian,
dan hukum distributif berlaku. Petunjuk: Invers multiplikatif dari $(a,b)$
adalah $\left(\frac{a}{a^2+b^2},\frac{-b}{a^2+b^2}\right)$.
\end{exercise}
\end{exbox}
